\chapter{Уравнение теплопроводности и индивидуальные границы}
\section{Тепловой баланс и закон Ньютона}
Пусть $S$ --- площадь поперечного сечения, $P$ --- его периметр,
$k>0$ --- теплопроводность, $c>0$ --- объёмная теплоёмкость,
$H\ge0$ --- коэффициент теплоотдачи, $F(x,t)$ --- объёмная мощность
источников, $T_e(x,t)$ --- температура среды. Для тонкого однородного
стержня температура $u$ приближённо постоянна по сечению.
Закон Фурье даёт поток вдоль оси $Q_x=-kS u_x$.
Приток через торцы малого элемента равен $kS u_{xx}\,dx$;
боковая потеря по закону Ньютона равна $HP(u-T_e)\,dx$.
Накопление тепла равно $cS u_t\,dx$. Поэтому
\[
 cS u_t=kS u_{xx}-HP(u-T_e)+FS,
\qquad u_t=a^2u_{xx}-bu+f,
\]
где
\[
 a^2=\frac{k}{c},\quad b=\frac{HP}{cS},\quad f=\frac Fc+bT_e.
\]
Для круглого сечения $S=\pi R^2$, $P=2\pi R$,
$b=2H/(cR)$. Это соответствует лекции 22, формуле (18),
с заменой лекционного $c_{\rm уд}\rho$ на $c$.
При теплоизоляции боковой поверхности $H=b=0$;
при отсутствии источников и бокового теплообмена также $f=0$.
\section{Условия варианта 11}
В столбцах 1--2 строки 11 указаны первый и второй род соответственно.
Поэтому произвольные индивидуальные данные записываем как
\[
 \boxed{u(0,t)=\mu(t),\qquad -u_x(l,t)=\nu(t),\qquad u(x,0)=\phi(x).}
\]
Знак справа сохраняет нормировку задания
$\gamma u(l,t)-\delta u_x(l,t)=\nu(t)$ при $\gamma=0$, $\delta=1$.
Это не означает, что $\nu$ уже является физическим потоком:
наружная плотность потока на правом торце равна $-ku_x(l,t)=k\nu(t)$.
При ином соглашении $u_x(l,t)=q(t)$ надо положить $\nu=-q$.
Для практического пункта 3 строка варианта задаёт
\[
 \mu(t)=A\cos(\omega_1t)+B\sin(\omega_2t),\qquad \nu(t)=0.
\]
В пунктах 2 и 4 по отдельному требованию задания вместо этих функций
берутся постоянные неоднородности $\mu=A$, $\nu=B$.
\section{Замечание о третьем роде}
Физический закон Ньютона на границе имеет вид
$k\partial_nu+H(u-T_e)=0$, где $n$ --- внешняя нормаль.
Следовательно, слева это $-ku_x+H(u-T_e)=0$,
справа --- $ku_x+H(u-T_e)=0$.
Запись с минусом перед $u_x$ на обоих концах и положительными
коэффициентами не является общей физической записью охлаждения обоих
торцов. В варианте 11 третьего рода нет, поэтому эту неоднозначность
не переносим в индивидуальное решение и не предполагаем положительность
спектра произвольной задачи с такими знаками.
\section{Согласование данных}
Для классического решения вплоть до $t=0$ необходимы
$\phi(0)=\mu(0)$ и $-\phi'(l)=\nu(0)$.
Без согласования рассматриваем решение при $t>0$, с начальным условием
в $L^2(0,l)$; непрерывность всех производных в угловых точках не заявляется.
